%! TEX root = master.tex
\chapter{4}{Parametric Curves}

A parametric curve gives $x$ and $y$ separately as functions of a parameter
$t$. This records not just the set of points but the way the curve is
traced: where it starts, which way it goes and how often it passes each
point. Circles are the basic building block. The chapter ends with the
geometry behind the rotary engine.

\section{Circles and arcs}\label{sec:parametric-circles}

%! TEX root = ../master.tex
\begin{theorem}[Parametric circle]\label{thm:parametric-circle}
  For $0 \leq t \leq 1$,
  \[
    x(t) = h + r \cos 2\pi t, \qquad y(t) = k + r \sin 2\pi t
  \]
  traces the circle with centre $(h, k)$ and radius $r$ exactly once,
  anticlockwise, starting and ending at $(h + r, k)$. Replacing $2\pi t$ by
  $\theta_0 + (\theta_1 - \theta_0)t$ traces only the arc from angle $\theta_0$
  to angle $\theta_1$.
\end{theorem}


\begin{exercise}\label{ex:parabola-path}
  Find parametric equations, with $0 \leq t \leq 1$, for a particle that
  moves along $y = x^2$ from $(1, 1)$ to $(-1, 1)$ and back to $(1, 1)$.
\end{exercise}

\begin{solution}
  Take
  \[
    x(t) = \cos 2\pi t, \qquad y(t) = \cos^2 2\pi t, \qquad 0 \leq t \leq 1.
  \]
  Then $y = x^2$ for every $t$, so the particle stays on the parabola. As $t$
  runs over $[0, \frac12]$, $2\pi t$ runs over $[0, \pi]$ and $x = \cos 2\pi t$
  decreases from $1$ to $-1$. As $t$ runs over $[\frac12, 1]$, $x$ increases
  back to $1$. So the particle goes from $(1, 1)$ to $(-1, 1)$ and returns,
  passing through each point of the arc once in each direction.
\end{solution}

\begin{exercise}\label{ex:smiley}
  Find parametric equations, with $0 \leq t \leq 1$ (several sets are
  needed), that draw the face in \cref{fig:smiley}.
\end{exercise}

\begin{figure}[htbp]
  \centering
  \includegraphics[width=0.4\linewidth]{parametric-face}
  \caption{The face, drawn from the four sets of equations below.}
  \label{fig:smiley}
\end{figure}

\begin{solution}
  All four pieces come from \cref{thm:parametric-circle}.
  \begin{enumerate}[label=(\alph*)]
    \item \emph{Face}, centre $(2, 2)$ and radius $2$:
      $x = 2 + 2\cos 2\pi t$, $y = 2 + 2\sin 2\pi t$.
    \item \emph{Left eye}, centre $(1, 3)$ and radius $0.2$:
      $x = 1 + 0.2\cos 2\pi t$, $y = 3 + 0.2\sin 2\pi t$.
    \item \emph{Right eye}, centre $(3, 3)$ and radius $0.2$:
      $x = 3 + 0.2\cos 2\pi t$, $y = 3 + 0.2\sin 2\pi t$.
    \item \emph{Smile}, the lower half of the circle with centre $(2, 2)$ and
      radius $1$, so the angles run from $\pi$ to $2\pi$:
      \[
        x = 2 + \cos(\pi + \pi t), \qquad y = 2 + \sin(\pi + \pi t).
      \]
      It starts at $(1, 2)$ when $t = 0$, passes through $(2, 1)$ when
      $t = \frac12$ and ends at $(3, 2)$ when $t = 1$. \qedhere
  \end{enumerate}
\end{solution}

\begin{exercise}\label{ex:two-circle-ellipse}
  Circles of radius $a$ and $b$, with $a > b > 0$, are centred at the origin.
  A ray from the origin at angle $t$ to the positive $x$-axis meets the
  large circle and then the small one. The point $P$ takes its
  $x$-coordinate from the first intersection and its $y$-coordinate from the
  second. Find parametric equations for the curve traced by $P$, and
  identify it.
\end{exercise}

\begin{figure}[htbp]
  \centering
  \includegraphics[width=0.5\linewidth]{two-circle-ellipse}
  \caption{$P$ takes its $x$-coordinate from the large circle and its
    $y$-coordinate from the small one.}
  \label{fig:two-circle-ellipse}
\end{figure}

\begin{solution}
  The ray at angle $t$ meets the circle of radius $a$ at
  $(a\cos t, a\sin t)$, and it meets the circle of radius $b$ at
  $(b\cos t, b\sin t)$.
  Hence
  \[
    x = a\cos t, \qquad y = b\sin t, \qquad 0 \leq t \leq 2\pi.
  \]
  Since $\cos^2 t + \sin^2 t = 1$, eliminating $t$ gives
  \[
    \frac{x^2}{a^2} + \frac{y^2}{b^2} = 1,
  \]
  which is the ellipse of \cref{thm:ellipse-standard} with semi-axes $a$ and
  $b$. As $t$ runs over $[0, 2\pi]$, $P$ goes once around it, anticlockwise.
\end{solution}

\section{The rotary engine}\label{sec:rotary-engine}

%! TEX root = ../master.tex
\begin{theorem}[Centroid]\label{thm:centroid}
  The three medians of the triangle with vertices $(x_i, y_i)$, $i = 1, 2, 3$,
  meet at the centroid
  \[
    \left( \frac{x_1 + x_2 + x_3}{3}, \frac{y_1 + y_2 + y_3}{3} \right),
  \]
  which lies two-thirds of the way along each median from its vertex.
\end{theorem}


\begin{exercise}\label{ex:epitrochoid}
  Show that an equilateral triangle can be inscribed in the epitrochoid
  \[
    x = 3r\cos t + b\cos 3t, \qquad y = 3r\sin t + b\sin 3t,
  \]
  and that as the triangle moves its centroid traces a circle of radius $b$
  about the centre of the epitrochoid.
\end{exercise}

\begin{proof}
  This is the principle of the rotary (Wankel) engine: a triangular rotor
  turns with its three tips on an epitrochoid-shaped housing, while its
  centre moves on a small circle.

  \emph{Vertices.} Fix $t$ and take the three points of the curve at the
  parameters $t_k = t + \frac{2\pi k}{3}$, for $k = 0, 1, 2$. Since
  $3t_k = 3t + 2\pi k$, the $b$-terms are the same for all three points:
  $b\cos 3t_k = b\cos 3t$ and $b\sin 3t_k = b\sin 3t$. Writing
  $G = (b\cos 3t, b\sin 3t)$, the vertices are
  \[
    V_k = G + 3r\,(\cos t_k, \sin t_k), \qquad k = 0, 1, 2.
  \]
  \emph{Equilateral.} The points $3r(\cos t_k, \sin t_k)$ lie on the circle
  of radius $3r$ about the origin, spaced $\frac{2\pi}{3}$ apart, so they form
  an equilateral triangle with side $2 \cdot 3r \sin\frac{\pi}{3} = 3\sqrt3\,r$.
  Adding the fixed point $G$ to all three is a translation, which preserves
  distances. So $V_0V_1V_2$ is equilateral, and its vertices lie on the
  epitrochoid.

  \emph{Centroid.} By \cref{thm:centroid}, the centroid is
  $G + r\sum_k (\cos t_k, \sin t_k)$. By the addition formulas, with
  $\cos\frac{2\pi}{3} = \cos\frac{4\pi}{3} = -\frac12$ and
  $\sin\frac{2\pi}{3} = -\sin\frac{4\pi}{3} = \frac{\sqrt3}{2}$,
  \[
    \sum_{k=0}^{2} \cos t_k
    = \cos t
      + \Bigl(-\tfrac12\cos t - \tfrac{\sqrt3}{2}\sin t\Bigr)
      + \Bigl(-\tfrac12\cos t + \tfrac{\sqrt3}{2}\sin t\Bigr) = 0,
  \]
  and likewise $\sum_k \sin t_k = 0$. So the centroid is
  \[
    G = (b\cos 3t,\; b\sin 3t),
  \]
  which, as $t$ varies, runs around the circle of radius $b$ centred at the
  origin.
\end{proof}

\begin{remark}
  The vertices turn at rate $1$ in $t$, while the centroid goes round at
  rate $3$. So the rotor's centre makes three orbits for every turn of the
  rotor, which is the 3:1 gearing of the Wankel engine.
\end{remark}
